\usection{Bases Teóricas}

\usubsection{Pensamiento computacional}

Según \cite{aho2011}, el pensamiento computacional es el conjunto de procesos del pensamiento involucrados en formular problemas tal que sus soluciones puedan ser representadas como pasos computacionales y algoritmos.

Por otro lado, \cite{wing2006} resalta un par de características para el pensamiento computacional, como que por ejemplo el pensamiento computacional se conceptualiza, no se programa, dando a entender que lo importante de este es la forma en la que abstrae un problema más que en la elaboración de su solución. Otro punto que resalta es que en el pensamiento computacional resaltan más las habilidades de creatividad e innovación que las de memorización. Por último resalta que el pensamiento computacional no es una forma de pensar como una computadora, sino más bien una habilidad humana fundamental que permite usar la lógica de la informática para resolver problemas en cualquier ámbito de la vida.

\usubsection{Lenguaje de programación}

Un lenguaje de programación se puede definir como un sistema de notación que se adapta al pensamiento humano, a través de él los programadores pueden desarrollar conjunto de órdenes e instrucciones para que un computador las ejecute \cite{trigo2004}. A diferencia de los lenguajes de máquina, que exigen un conocimiento profundo del código binario, estos lenguajes resultan muy útiles, ya que actúan como un puente que permite a los programadores desarrollar software sin tener que interactuar directamente con el lenguaje de la computadora.

\usubsection{Tesis de Church-Turing}

\cite{aho2011} explica que la teoría de la computación explora las capacidades y limitaciones de los modelos de computación. El modelo más importante para la computación secuencial es la Máquina de Turing, propuesta por Alan Turing en 1936. La tesis de Church-Turing postula que cualquier función que pueda ser computada, puede serlo por una Máquina de Turing. Esto significa que modelos de computación como el cálculo lambda, las máquinas de acceso aleatorio y todos los lenguajes de programación de propósito general son equivalentes a las Máquinas de Turing, ya que pueden computar el mismo conjunto de funciones. Aho destaca que, aunque se utilizan diferentes notaciones y lenguajes de programación por ser más fáciles de escribir y entender, la tesis de Church-Turing sugiere que no importa la notación que se elija, una Máquina de Turing puede computar cualquier función.

\usubsection{Lenguaje de programación Turing-completo}

Un lenguaje de programación se considera Turing-completo cuando es capaz de realizar cualquier cómputo que una máquina de Turing teórica puede llevar a cabo. Para lograr esto, debe incluir tres elementos esenciales: operaciones secuenciales para ejecutar instrucciones una tras otra, operaciones condicionales que permiten tomar decisiones lógicas, e iteración o recursión ilimitada, lo que posibilita la repetición de un proceso sin un número fijo de veces. La combinación de estos componentes básicos garantiza que el lenguaje tenga el poder de simular la lógica de cualquier algoritmo, lo que lo hace universalmente aplicable para resolver una amplia gama de problemas computacionales \cite{zeil2024}.

\usubsubsection{Programación por flujos de datos (dataflow)}

\cite{tran2024} indica que la programación por flujos de datos es un paradigma de programación en el que se modelan programas como grafos dirigidos consistentes de nodos que representan operaciones que aceptan una entrada y emiten una salida hacia el o los siguientes nodos. Los programas orientados a dataflow describen las transformaciones de datos con bloques ya establecidos y fáciles de entender, además que simplifica la ejecución paralela del algoritmo, simplemente visualizando las dependencias de datos en los distintos nodos. Según \cite{hasselbring2021}, la programación en dataflow se preocupa de hacia dónde los datos son transportados a través del sistema, y qué transformaciones son aplicadas durante este proceso. Esto contrasta con el control flow, donde la preocupación principal es el orden de las operaciones.

\usubsection{Concurrencia}

Según \cite{marinescu2013}, la concurrencia implica la ejecución simultánea de varias actividades, permitiendo reducir el tiempo de ejecución de problemas con gran carga de datos. Para aprovechar la concurrencia es necesario diseñar algoritmos paralelos. En el diseño de sistemas, la concurrencia puede utilizarse con distintos objetivos como ocultar la latencia, mejorar el rendimiento y evitar limitaciones de recursos físicos. La concurrencia se ha vuelto algo ubicuo, disponible en sistemas embebidos, aplicaciones paralelas y aplicaciones distribuidas, donde habilita la elasticidad en entornos de computación en la nube.
